% ===========================================================================
\part{The mental model}
% ===========================================================================

\chapter{One source tree, two programs}
\label{ch:twobuilds}

\section{The fact everything else follows from}

\begin{keyidea}
\file{src/lib.rs} is compiled twice, into two different programs, for two different machines. Every
module, every struct and every dependency is either in both programs, or in exactly one of them,
and it is your job --- not the compiler's --- to say which.
\end{keyidea}

This is the single idea that, once internalised, makes Leptos stop feeling arbitrary. Read
\file{Cargo.toml} again with it in mind:

\begin{lstlisting}[style=toml,caption={The features that select which program is being built.}]
[lib]
crate-type = ["cdylib", "rlib"]   # cdylib -> wasm target; rlib -> so main.rs can `use` it

[features]
ssr     = ["dep:actix-web", "dep:actix-files", "dep:leptos_actix", "dep:diesel",
           "dep:serde", "dep:serde_json",
           "leptos/ssr", "leptos_meta/ssr", "leptos_router/ssr"]
hydrate = ["leptos/hydrate"]
csr     = ["leptos/csr"]
\end{lstlisting}

\code{ssr} and \code{hydrate} are not two configurations of one program. They are the names of the
two programs. Figure~\ref{fig:twobuilds} shows what \code{cargo leptos build} actually does with
them.

\begin{diagram}{What \code{cargo leptos build} produces. The same \file{src/lib.rs} is read by both
compilations; the feature flag decides which parts of it exist.}{fig:twobuilds}
\begin{tikzpicture}[node distance=5mm]
  \node[shared,minimum width=32mm] (src) {\textbf{src/lib.rs}\\[1pt]\scriptsize + everything it declares};

  % left branch: server
  \node[srv,below left=14mm and 12mm of src,minimum width=46mm] (sb)
    {\code{cargo build --bin egonik-site}\\[1pt]\scriptsize
     \code{--features ssr --no-default-features}\\\scriptsize target: x86\_64 native};
  \node[srv,below=7mm of sb,minimum width=46mm] (sout)
    {\textbf{target/debug/egonik-site}\\[1pt]\scriptsize an executable: Actix + Diesel + your domain};

  % right branch: client
  \node[cli,below right=14mm and 12mm of src,minimum width=46mm] (cb)
    {\code{cargo build --lib}\\[1pt]\scriptsize
     \code{--features hydrate --no-default-features}\\\scriptsize target: wasm32-unknown-unknown};
  \node[cli,below=7mm of cb,minimum width=46mm] (cbg)
    {\code{wasm-bindgen}\\[1pt]\scriptsize generates the JS glue};
  \node[cli,below=7mm of cbg,minimum width=46mm] (cout)
    {\textbf{target/site/pkg/egonik-site.\{wasm,js\}}\\[1pt]\scriptsize downloaded by the browser};

  \draw[flow] (src.south) -- ++(0,-4mm) -| node[lbl,pos=0.25,above]{\code{ssr}} (sb.north);
  \draw[flow] (src.south) -- ++(0,-4mm) -| node[lbl,pos=0.25,above]{\code{hydrate}} (cb.north);
  \draw[flow] (sb) -- (sout);
  \draw[flow] (cb) -- (cbg);
  \draw[flow] (cbg) -- (cout);

  \node[box,fill=wash,below=9mm of sout,minimum width=46mm,align=center] (run)
    {\textbf{the running server}\\[1pt]\scriptsize serves the wasm bundle it was built alongside};
  \draw[flow] (sout) -- (run);
  \draw[flow,dashed] (cout.south) |- ++(0,-8mm) -| (run.east);
\end{tikzpicture}
\end{diagram}

Two consequences are worth stating explicitly, because both bite in practice.

\textbf{First: a dependency you add is, by default, in both programs.} It is not enough to only
\emph{use} a crate on the server; if it is listed unconditionally under \code{[dependencies]}, it
is compiled into the browser bundle too. That is measurable in the current code.

\begin{verified}[Verified: what is actually in the wasm bundle today]
\begin{lstlisting}[style=out]
$ cargo tree --target wasm32-unknown-unknown --no-default-features \
             --features hydrate --depth 1
egonik-site v0.1.0
|-- anyhow v1.0.104
|-- chrono v0.4.45
|-- console_error_panic_hook v0.1.7
|-- dotenvy v0.15.7          <-- reads .env files. In a browser.
|-- leptos v0.8.20
|-- leptos_meta v0.8.6
|-- leptos_router v0.8.15
`-- wasm-bindgen v0.2.126
\end{lstlisting}
\end{verified}

\code{dotenvy} in a WebAssembly bundle is not merely wasteful, it is nonsensical: it exists to read
a file from a filesystem the browser does not have. \code{anyhow} is server-side error plumbing.
\code{chrono} is genuinely needed on both sides if dates are rendered --- but only its
serialisation and formatting parts, not its system-clock parts. These are defects
\dnum{9} and \dnum{10}.

\textbf{Second: \code{\#[cfg(feature = "ssr")]} is a scalpel, and it is currently being used as a
tourniquet.} Look at what \file{src/lib.rs} does:

\begin{lstlisting}[style=rs]
#[cfg(feature = "ssr")] pub mod core;
#[cfg(feature = "ssr")] pub mod database;
#[cfg(feature = "ssr")] pub mod job_experience;
#[cfg(feature = "ssr")] pub mod personal_information;
#[cfg(feature = "ssr")] pub mod portfolio;
#[cfg(feature = "ssr")] pub mod publications;
#[cfg(feature = "ssr")] pub mod schema;
\end{lstlisting}

Gating \code{database} and \code{schema} is correct and necessary --- Diesel must never reach the
browser. But gating the entire \code{portfolio}, \code{publications}, \code{job\_experience} and
\code{personal\_information} modules means the \emph{types describing your content} exist only on the
server. Chapter~\ref{ch:boundary} shows why that becomes a wall the moment you write your first
server function, and Chapter~\ref{ch:dto} gives the split that removes it.

\begin{why}[Why this gate was added, and why that was reasonable]
The gate is almost certainly there because without it the wasm build failed. It \emph{would} have:
the models derive \code{Serialize}, \code{Queryable} and \code{Selectable}, and none of those
derive macros are available in the \code{hydrate} build. Adding the gate made the error go away.
The problem is that it made the error go away by removing the types from the client entirely,
which is the opposite of what you will need. The right fix separates the Diesel-flavoured struct
from the wire-flavoured struct, rather than hiding both.
\end{why}

\section{Reading a \code{cfg} error correctly}

When either build breaks, the first diagnostic question is always \emph{which of the two programs
is failing?} \code{cargo leptos} interleaves both, so the output can be confusing. Check them
independently:

\begin{lstlisting}[style=sh]
# program 1: the server
cargo check --no-default-features --features ssr

# program 2: the browser
cargo check --lib --no-default-features --features hydrate \
            --target wasm32-unknown-unknown
\end{lstlisting}

\begin{verified}[Verified: both programs compile today]
Both commands complete successfully against the current tree --- the server in 18.5\,s, the wasm
library in 59.3\,s from a warm cache. Keep it that way; make these two commands the definition of
``it builds''.
\end{verified}

\begin{pattern}[Make the two-build check a habit]
Before every commit, run both commands. A change that satisfies \code{cargo leptos watch} can
still have broken the other target, because \code{watch} may not have rebuilt it yet. This is also
the single most valuable thing to put in CI --- see Chapter~\ref{ch:ops}.
\end{pattern}

\section{A note on the \code{csr} feature}

\file{Cargo.toml} also defines \code{csr}, and \file{src/main.rs} carries two extra \code{main}
functions to support it. Client-side rendering means shipping an empty HTML shell and letting the
wasm bundle build the whole page --- no server rendering at all. It is useful for testing
components under \code{trunk serve}, and useless for a personal website, where server rendering is
the entire point (search engines, first-paint speed, working without JavaScript).

You have three configurations to keep straight and only two that you use. Deleting \code{csr}
from \file{Cargo.toml} and removing the two dead \code{main} functions from \file{src/main.rs}
costs ten minutes and permanently shrinks the space of states you have to reason about. This is
recorded as defect \dnum{20}, at low severity, because it is housekeeping rather than a fault.

% ===========================================================================
\chapter{The life of a request}
\label{ch:lifecycle}

Understanding \emph{when} your code runs is what allows you to place it correctly. A Leptos
component function may execute on the server, in the browser, or both, and it will do different
things in each case. This chapter follows four distinct journeys through the system.

\section{Journey 1: a cold page load}

Someone types your URL. Nothing is cached. Figure~\ref{fig:cold} traces what happens.

\begin{diagram}{A cold page load. Note that your component function runs \emph{twice} --- once on the
server at step 4, once in the browser at step 9 --- and that the data fetched at step 5 is
serialised into the HTML so the browser does not have to fetch it again.}{fig:cold}
\begin{tikzpicture}[x=1mm,y=1mm]
  \def\lh{106}
  % lanes
  \foreach \x/\nm/\st in {0/Browser/cli, 34/Actix/srv, 68/Leptos (ssr)/srv, 100/Repository/srv, 130/Postgres/srv} {
    \node[\st,minimum width=22mm,anchor=north,inner sep=3pt] at (\x,0) {\textbf{\nm}};
    \draw[draw=rule,line width=0.6pt,dashed] (\x,-9) -- (\x,-\lh);
  }
  % steps
  \def\stp#1#2#3#4{\draw[flowa] (#1,#3) -- (#2,#3) node[lbl,midway,above=0.3mm]{#4};}
  \stp{0}{34}{-17}{1. \code{GET /}}
  \stp{34}{68}{-26}{2. no static match $\rightarrow$ \code{leptos\_routes}}
  \node[lbl,anchor=west,align=left] at (70,-34) {3. match the route in \code{<Routes>}};
  \node[lbl,anchor=west,align=left] at (70,-41) {4. \textbf{run your component} $\rightarrow$ it hits a \code{Resource}};
  \stp{68}{100}{-49}{5. server fn body runs \emph{in-process}}
  \stp{100}{130}{-57}{6. SQL}
  \draw[flow] (130,-64) -- (100,-64) node[lbl,midway,above=0.3mm]{rows};
  \draw[flow] (100,-71) -- (68,-71) node[lbl,midway,above=0.3mm]{7. \code{Vec<Dto>}};
  \node[lbl,anchor=west,align=left] at (70,-78) {8. render HTML \emph{and} embed the DTOs as JSON};
  \draw[flow] (68,-86) -- (34,-86);
  \draw[flow] (34,-86) -- (0,-86) node[lbl,midway,above=0.3mm]{HTML $+$ payload};
  \node[lbl,anchor=west,align=left] at (2,-93) {9. paint, then fetch the wasm, then \textbf{hydrate}};
  \node[lbl,anchor=west,align=left] at (2,-100) {10. the \code{Resource} reads the embedded JSON --- \emph{no second request}};
\end{tikzpicture}
\end{diagram}

Three details in that figure carry real weight.

\textbf{At step 5, the server function is not an HTTP call.} When a \code{\#[server]} function is
invoked during server rendering, Leptos calls the function body directly, in the same process, on
the same request. There is no localhost round trip. This is why a server function is cheap during
SSR and why it is safe to have several of them on one page.

\textbf{At step 8, the data is serialised into the HTML.} This is the mechanism that makes
hydration fast, and it is also the reason your DTOs must implement \code{Serialize} on the server
\emph{and} \code{Deserialize} in the browser. Not one or the other --- both, in the same type, in
both builds. Chapter~\ref{ch:boundary} is entirely about this constraint.

\textbf{At step 9, the page is already visible before the wasm arrives.} Server-rendered HTML is
real HTML. If the wasm bundle fails to download, the visitor still sees your content; they just
cannot interact with anything that requires reactivity. For a personal website that is a very good
failure mode, and it is worth preserving deliberately by keeping content in server-rendered markup
rather than fetching it client-side after load.

\section{Journey 2: navigating to a second page}

The visitor clicks a link to \code{/portfolio}. Now the wasm bundle is loaded and Leptos is in
control.

\begin{diagram}{Client-side navigation. Actix serves no HTML at all; only the server function is
called, and only the data crosses the network.}{fig:warm}
\begin{tikzpicture}[x=1mm,y=1mm]
  \foreach \x/\nm/\st in {0/Browser (wasm)/cli, 46/Actix/srv, 88/server fn body/srv, 126/Postgres/srv} {
    \node[\st,minimum width=26mm,anchor=north,inner sep=3pt] at (\x,0) {\textbf{\nm}};
    \draw[draw=rule,line width=0.6pt,dashed] (\x,-9) -- (\x,-62);
  }
  \node[lbl,anchor=west,align=left] at (2,-16) {1. the click is intercepted; \code{history.pushState}};
  \node[lbl,anchor=west,align=left] at (2,-23) {2. the new route's component runs \textbf{in the browser}};
  \draw[flowa] (0,-31) -- (46,-31) node[lbl,midway,above=0.3mm]{3. \code{POST} to the generated endpoint};
  \draw[flowa] (46,-38) -- (88,-38);
  \draw[flow] (88,-45) -- (126,-45) node[lbl,midway,above=0.3mm]{SQL};
  \draw[flow] (126,-52) -- (88,-52);
  \draw[flow] (88,-59) -- (0,-59) node[lbl,pos=0.6,above=0.3mm]{4. serialised \code{Vec<Dto>} --- no HTML at all};
\end{tikzpicture}
\end{diagram}

The same component function you wrote for journey 1 now runs in WebAssembly, and the same
\code{\#[server]} function you called now performs an HTTP request instead of a direct call. You
did not write two versions of anything. That is the payoff of the two-build model, and it is why
the discipline it demands is worth accepting.

\section{Journey 3: no JavaScript at all}

If the visitor has JavaScript disabled, or the bundle 404s, journey 1 still completes through
step 8. They get a fully rendered page. Forms submitted through Leptos's \code{<ActionForm>} even
continue to work, because they degrade to ordinary HTML form posts.

\begin{pattern}[Design for the no-wasm case first]
Render content on the server. Use reactivity for interaction, not for content delivery. A personal
website whose text only appears after 400\,KB of WebAssembly has loaded is strictly worse than one
whose text is in the HTML --- for readers, for search engines, and for you when you are debugging.
\end{pattern}

\section{Journey 4: a write}

You do not have writes yet, so this is the shape to aim for rather than a description of
existing behaviour. A write is a server function that mutates, wrapped in an \code{Action} rather
than a \code{Resource}:

\begin{itemize}
  \item a \code{Resource} is for \emph{reading}: it runs automatically when its inputs change, and
    its result is serialised into SSR output;
  \item an \code{Action} is for \emph{writing}: it runs when you tell it to, exactly once per
    dispatch, and never runs during server rendering.
\end{itemize}

Getting this pairing wrong is a classic Leptos error --- a mutation inside a \code{Resource} will
fire during server rendering and again on hydration, and you will insert two rows. Keep the rule:
\textbf{reads are Resources, writes are Actions.}

\section{Where the current code sits in these journeys}

Nowhere, yet, and that is the point of the chapter. \file{src/app.rs} has a \code{RwSignal}
counter, which participates only in journey 2, and only in the browser. There is no
\code{Resource}, so no data is fetched at step 5; no server function, so nothing is registered at
step 3 of journey 2; and no \code{provide\_context}, so nothing the repositories need would be
reachable if there were. Chapter~\ref{ch:composition} wires it, and
Chapter~\ref{ch:loading} writes the first \code{Resource}.

% ===========================================================================
\chapter{The boundary: what may enter the browser}
\label{ch:boundary}

\section{The rule}

\begin{keyidea}
A type that crosses between server and browser must compile in \emph{both} programs, must derive
\code{Serialize} and \code{Deserialize} in \emph{both} programs, and must not transitively depend
on anything server-only. A type that does not cross must be unavailable to the browser
\emph{on purpose}.
\end{keyidea}

Figure~\ref{fig:boundary} draws the boundary. The left column is server-only, the right is shared,
and the seam between them is where DTOs live.

\begin{diagram}{The boundary. Everything on the left is invisible to the browser by construction;
everything in the middle column is compiled into both programs and must therefore be free of
server-only dependencies.}{fig:boundary}
\begin{tikzpicture}[node distance=4mm]
  % server column
  \node[srv,minimum width=40mm] (r) {\code{diesel::table!} schema};
  \node[srv,minimum width=40mm,below=of r] (m) {Diesel models\\\scriptsize \code{Queryable}, \code{Selectable}};
  \node[srv,minimum width=40mm,below=of m] (rep) {repositories\\\scriptsize the pool, SQL, \code{r2d2}};
  \node[srv,minimum width=40mm,below=of rep] (dom) {domain rules, validation};
  \node[srv,minimum width=40mm,below=of dom] (cfg) {config, secrets, \code{dotenvy}};

  % shared column
  \node[shared,text width=34mm,right=24mm of m] (dto) {\textbf{DTOs}\\\scriptsize plain structs\\\scriptsize \code{Serialize + Deserialize}};
  \node[shared,text width=34mm,below=8mm of dto] (sig) {server-fn signatures\\\scriptsize \code{\#[server] fn}\\\scriptsize \code{-> Result<Dto>}};
  \node[shared,text width=34mm,below=8mm of sig] (err) {error enum\\\scriptsize the parts the UI shows};

  % client column
  \node[cli,text width=32mm,right=22mm of dto] (comp) {components, \code{view!}};
  \node[cli,text width=32mm,below=8mm of comp] (res) {\code{Resource}, \code{Action}};
  \node[cli,text width=32mm,below=8mm of res] (fmt) {formatting, \code{chrono} display};

  % seam: a vertical line midway between the server and shared columns,
  % spanning the full height of the server column plus a little margin
  \coordinate (sx)   at ($(m.east)!0.5!(dto.west)$);
  \coordinate (stop) at ($(sx |- r.north)   + (0, 8mm)$);
  \coordinate (sbot) at ($(sx |- cfg.south) + (0,-12mm)$);
  \draw[seam] (stop) -- (sbot)
    node[pos=0.45,fill=white,inner sep=2pt,rotate=90,
         font=\sffamily\scriptsize\bfseries,text=crimson]{the boundary};

  \draw[flow] (m) -- node[lbl,above]{\code{From}} (dto);
  \draw[flow] (dto) -- (comp);
  \draw[flow] (sig) -- (res);

  \begin{scope}[on background layer]
    \node[lane,fit=(r)(cfg),inner sep=3mm,label={[tag]below:server build only (\code{ssr})}] {};
    \node[lane,fit=(dto)(err),inner sep=3mm,label={[tag]below:both builds}] {};
    \node[lane,fit=(comp)(fmt),inner sep=3mm,label={[tag]below:runs in the browser}] {};
  \end{scope}
\end{tikzpicture}
\end{diagram}

\section{The failure you are going to hit}

This is not hypothetical, so it is worth showing exactly. Suppose you decide to return your
existing model straight from a server function --- the obvious first attempt:

\begin{lstlisting}[style=rs]
// the natural thing to write, and it cannot work
#[server]
pub async fn list_publications() -> Result<Vec<PublicationItem>, ServerFnError> { ... }
\end{lstlisting}

Two things break at once.

\textbf{The type does not exist in the browser build.} \code{PublicationItem} lives in
\file{src/publications/model.rs}, and \file{src/lib.rs} gates that module on \code{ssr}. The
\code{hydrate} build has never heard of it, and the \code{\#[server]} macro needs the signature to
compile on \emph{both} sides --- it generates the client stub from the same declaration.

\textbf{And \code{serde} is not available in the browser build either.} \code{serde} is declared
\code{optional = true} and enabled only by the \code{ssr} feature. Remove the module gate and you
trade one error for another:

\begin{verified}[Verified: the exact error, reproduced]
A struct deriving \code{serde::Serialize} was added to the crate root without a \code{cfg} gate,
and the client program was compiled:
\begin{lstlisting}[style=out]
$ cargo check --lib --no-default-features --features hydrate \
              --target wasm32-unknown-unknown
error[E0432]: unresolved import `serde`
error: could not compile `egonik-site` (lib) due to 1 previous error
\end{lstlisting}
The file was then removed and the tree restored. This is defect \dnum{6}.
\end{verified}

\begin{pitfall}[The confusing part]
\code{serde} \emph{is} physically compiled into the wasm build --- Leptos depends on it. But
because your crate declares it as an optional dependency that only \code{ssr} enables, \emph{your}
code is not permitted to name it. The crate is present and the path is unresolvable. That is why
the error message is \code{unresolved import} rather than something about a missing dependency,
and it is why the fix is in \file{Cargo.toml}, not in the \code{use} statement.
\end{pitfall}

\section{The fix, in two lines of \code{Cargo.toml}}

\begin{lstlisting}[style=toml,caption={\code{serde} becomes unconditional; the server-only crates stay optional.}]
[dependencies]
serde = { version = "1", features = ["derive"] }         # both programs need it
serde_json = { version = "1", optional = true }          # server only, unless you need it client-side
chrono = { version = "0.4", features = ["serde"], default-features = false }
diesel = { version = "2.3", optional = true, features = [...] }
dotenvy = { version = "0.15", optional = true }           # <-- make optional
anyhow  = { version = "1",    optional = true }           # <-- make optional

[features]
ssr = ["dep:actix-web", "dep:actix-files", "dep:leptos_actix", "dep:diesel",
       "dep:serde_json", "dep:dotenvy", "dep:anyhow",
       "leptos/ssr", "leptos_meta/ssr", "leptos_router/ssr"]
\end{lstlisting}

\code{serde} is a small crate and is already in the bundle; making it unconditional costs nothing
and removes an entire category of error. \code{dotenvy} and \code{anyhow} move the other way ---
into \code{ssr} --- which is defect \dnum{9}'s fix and shrinks the wasm bundle.

Note the \code{chrono} line. Dates \emph{do} need to cross the boundary --- a publication year, an
employment period --- so \code{chrono} must stay unconditional, but with
\code{default-features = false} it drops the parts that want a system clock and a timezone
database, which is most of its weight. Keep \code{features = ["serde"]} so \code{NaiveDate}
serialises.

\section{The checklist}

Before adding any type or dependency, ask in order:

\begin{enumerate}
  \item \textbf{Does this need to be visible to a component?} If no, gate it on \code{ssr} and
    stop. Most things stop here: queries, the pool, migrations, secrets.
  \item \textbf{Does it need to travel over the wire?} If yes, it is a DTO. It goes in the
    ungated \file{src/dto/} module, derives \code{Serialize + Deserialize}, and contains
    \emph{only} primitives, \code{String}, \code{Option}, \code{Vec} and other DTOs --- plus
    \code{NaiveDate}, which is fine.
  \item \textbf{Does it pull in a crate?} Check the crate compiles for
    \code{wasm32-unknown-unknown} and does not want the filesystem, threads, or a socket. When in
    doubt, run the two-build check from Chapter~\ref{ch:twobuilds}.
\end{enumerate}

\begin{antipattern}[Reaching for \code{cfg} to silence an error]
When the wasm build complains about a type, the tempting fix is \code{\#[cfg(feature = "ssr")]}.
Sometimes that is right --- but if the type is something a page needs to display, the \code{cfg}
is postponing the problem, not solving it. The question to ask is not ``how do I hide this from
the client'' but ``does the client need this shape, and if so what should that shape be?'' The
answer to the second question is a DTO, which is Chapter~\ref{ch:dto}.
\end{antipattern}

\section{A per-dependency verdict for this \file{Cargo.toml}}

The checklist above is general. Here it is applied to every dependency currently declared, because
four of the eleven are on the wrong side and two are not used at all.

\begin{table}[htbp]
\centering
\caption{Every declared dependency, and where it belongs.}
\label{tab:deps}
\small
\begin{tabular}{@{}L{0.20\textwidth}L{0.15\textwidth}L{0.15\textwidth}L{0.42\textwidth}@{}}
\toprule
\textbf{Crate} & \textbf{Declared} & \textbf{Should be} & \textbf{Why} \\
\midrule
\code{leptos}, \code{leptos\_meta}, \code{leptos\_router} & both & \textcolor{moss}{both} & the
framework; each has its own \code{ssr} feature \\
\code{actix-web}, \code{actix-files}, \code{leptos\_actix} & \code{ssr} & \textcolor{moss}{\code{ssr}} & correct \\
\code{diesel} & \code{ssr} & \textcolor{moss}{\code{ssr}} & correct, and mandatory: the
\code{postgres} feature links native \code{libpq} \\
\code{serde} & \textbf{\code{ssr}} & \textcolor{crimson}{\textbf{both}} & \dnum{6}. Wire types must
serialise in the browser \\
\code{serde\_json} & \code{ssr} & \textcolor{crimson}{delete} & used nowhere in \file{src/} \\
\code{http} & optional, \emph{no feature} & \textcolor{crimson}{delete} & enabled by nothing,
used by nothing \\
\code{chrono} & both & \textcolor{amber}{both, trimmed} & \dnum{10}. Dates do cross; add
\code{default-features = false} \\
\code{dotenvy} & both & \textcolor{crimson}{\code{ssr}} & \dnum{9}. Reads \file{.env} from a
filesystem the browser has not got \\
\code{anyhow} & both & \textcolor{crimson}{\code{ssr}} & \dnum{9}. Server-side error plumbing \\
\code{wasm-bindgen} & both & \textcolor{amber}{\code{hydrate}} & needed only by
\code{pub fn hydrate()}; currently also compiled natively \\
\code{console\_error\_panic\_hook} & both & \textcolor{amber}{\code{hydrate}} & browser-only ---
see the next chapter for why it matters \\
\bottomrule
\end{tabular}
\end{table}

\begin{verified}[Verified: two of these are entirely dead]
\begin{lstlisting}[style=out]
$ grep -n 'http' Cargo.toml
13:http = { version = "1.3.1", optional = true }     <-- in no feature list
$ grep -rn 'http::' src/
src/app.rs:60: resp.set_status(actix_web::http::StatusCode::NOT_FOUND);   <-- actix's, not this one
$ grep -rn 'serde_json' src/
  (no output)
\end{lstlisting}
Both declarations can go. \code{http} in particular is invisible: because no feature enables it,
it is never compiled, so it produces no warning and no build cost --- just a line that implies a
dependency you do not have. These are defects \dnum{34} and \dnum{35}.
\end{verified}

% ===========================================================================
\chapter{The run contract, and how to tell hydration happened}
\label{ch:runcontract}

\section{Why the binary behaves differently from \code{cargo leptos watch}}

\begin{keyidea}
The compiled server binary is not self-contained. It reads its own configuration from
\code{LEPTOS\_*} environment variables at startup, and \code{cargo leptos} sets them for you. Run
the binary yourself without them and you get a site that renders perfectly and never becomes
interactive.
\end{keyidea}

This is the single most confusing failure available in this stack, because every visible signal says
the application is fine. It is worth showing exactly, since you will meet it the first time you run
the binary directly or put it in a container.

\begin{verified}[Verified: running \file{target/debug/egonik-site} with no \code{LEPTOS\_*} variables]
The binary prints its own warning, and then serves a page:
\begin{lstlisting}[style=out]
$ env -u LEPTOS_OUTPUT_NAME -u LEPTOS_SITE_ROOT ./target/debug/egonik-site
It looks like you're trying to compile Leptos without the LEPTOS_OUTPUT_NAME environment
variable being set. There are two options
 1. cargo-leptos is not being used, but get_configuration() is being passed None. This
    needs to be changed to Some("Cargo.toml")
 2. You are compiling Leptos without LEPTOS_OUTPUT_NAME being set with cargo-leptos.
    This shouldn't be possible!
listening on http://127.0.0.1:3000
\end{lstlisting}

And here is what that page contains:
\begin{lstlisting}[style=out]
$ curl -s localhost:3000/ | grep -oE '(src|href)="/pkg/[^"]*"'
href="/pkg/_bg.wasm"                 <-- the output name is EMPTY
href="/pkg/egonik-site.css"          <-- but this one is hardcoded in app.rs, so it works
href="/pkg/.js"                      <-- and this one is broken

$ curl -s -o /dev/null -w '%{http_code}' localhost:3000/pkg/.js
400

$ curl -s localhost:3000/ | grep -o 'Click Me[^<]*'
Click Me:                            <-- the HTML is complete and correct
\end{lstlisting}
\end{verified}

Read that carefully, because the shape of the failure is the lesson. The server-rendered HTML is
complete. The stylesheet loads --- because \file{app.rs} hardcodes
\code{href="/pkg/egonik-site.css"} rather than deriving it. So the page \emph{looks entirely
correct} in a browser and in view-source. Only the WebAssembly is missing, so nothing is
interactive: the counter button does nothing, and no error appears anywhere except a 400 in the
network tab for a file called \code{.js}.

\begin{pitfall}[The two options in that warning message, and which one you want]
The message offers two remedies. Take neither at face value.

Option 1 --- ``pass \code{Some("Cargo.toml")} to \code{get\_configuration}'' --- makes the running
binary read \file{Cargo.toml} from disk at startup. It works in your source tree and fails in any
container that ships only the binary and the site directory. Do not do this.

The right answer is to \textbf{set the environment variables}, which is what \code{cargo leptos}
does in development and what your \file{Dockerfile} must do in production:
\begin{lstlisting}[style=sh]
LEPTOS_OUTPUT_NAME=egonik-site \
LEPTOS_SITE_ROOT=target/site \
LEPTOS_SITE_PKG_DIR=pkg \
LEPTOS_ENV=PROD \
./target/release/egonik-site
\end{lstlisting}
With \code{LEPTOS\_OUTPUT\_NAME} set, the link becomes \code{/pkg/egonik-site.js} and returns 200.
\end{pitfall}

\begin{table}[htbp]
\centering
\caption{Symptoms of a missing runtime variable. None of these produce a Rust error.}
\label{tab:runtimeenv}
\small
\begin{tabular}{@{}L{0.25\textwidth}L{0.30\textwidth}L{0.38\textwidth}@{}}
\toprule
\textbf{Variable} & \textbf{If unset} & \textbf{Symptom} \\
\midrule
\code{LEPTOS\_OUTPUT\_NAME} & bundle name is empty & \code{/pkg/.js} $\rightarrow$ 400; page
renders, never hydrates \\
\code{LEPTOS\_SITE\_ROOT} & defaults to \code{"."} & static files 404; page renders unstyled \\
\code{LEPTOS\_SITE\_ADDR} & \code{127.0.0.1:3000} & in a container: unreachable from outside.
Set \code{0.0.0.0:8080} \\
\code{LEPTOS\_ENV} & \code{DEV} & the live-reload client is injected and retries
\code{localhost:3001} forever \\
\bottomrule
\end{tabular}
\end{table}

\section{How to tell whether hydration actually happened}

Because a non-hydrating page looks identical to a working one, you need a deliberate test. Three
checks, cheapest first:

\begin{enumerate}
  \item \textbf{Does \code{/pkg/<output-name>.js} return 200?} If not, nothing else matters. This is
    the check that catches the entire failure above, and it is one \code{curl}.
  \item \textbf{Does an interactive element work?} The starter template's counter is genuinely
    useful for this, which is a reason to keep something clickable on the page until you have real
    interaction. If the button increments, the wasm loaded and took over.
  \item \textbf{Is the browser console clean?} Which brings us to the panic hook.
\end{enumerate}

\begin{why}[What \code{console\_error\_panic\_hook} is actually for]
\file{src/lib.rs} calls \code{console\_error\_panic\_hook::set\_once()} inside \code{hydrate()}.
Without it, a panic in WebAssembly aborts with the message \code{unreachable executed} and no
indication of where. With it, you get the real panic message and a stack trace in the browser
console.

This is not optional infrastructure --- it is the only diagnostic you have on the client side.
Keep the call, and make ``check the console'' part of how you verify a page, because a hydration
panic silently leaves the page in the same non-interactive state as a missing bundle.
\end{why}

\section{The component body runs twice, and that has consequences}

From Chapter~\ref{ch:lifecycle}: on a cold load your component function executes on the server, and
then again in the browser during hydration. Leptos requires that both executions produce the
\emph{same} markup. If they differ, hydration mismatches --- the framework attaches event listeners
to the wrong nodes, and behaviour becomes inexplicable.

\begin{antipattern}[Anything non-deterministic inside a \code{view!}]
\begin{lstlisting}[style=rs]
// renders one time on the server, a different time in the browser
view! { <p>"Rendered at " {chrono::Local::now().to_string()}</p> }

// a different value on each of the two runs
view! { <li id=format!("item-{}", rand::random::<u32>())>...</li> }
\end{lstlisting}
Both compile. Both produce a mismatch. Anything time-dependent or random must either be computed
once on the server and passed through as data, or computed only on the client inside an
\code{Effect}, which does not run during server rendering.
\end{antipattern}

\begin{pitfall}[\code{chrono::Local} needs a timezone the browser will not agree about]
This is a live risk for this project, because \code{chrono} is compiled into both builds and the
schema is full of dates. \code{Local::now()} resolves against the machine's timezone --- the
server's timezone on one run, the visitor's on the other. Format dates on the server, or format
them from a \code{NaiveDate} with no timezone involved at all, which is what every date column in
this schema already is.
\end{pitfall}

\section{Where the assets actually end up}

One more piece of the run contract, and a genuinely surprising one. \file{Cargo.toml} says
\code{assets-dir = "assets"}. The contents of that directory are copied \emph{into the root of the
site directory}, not into a subdirectory named \code{assets}.

\begin{verified}
\begin{lstlisting}[style=out]
$ find target/site -maxdepth 2
target/site
target/site/favicon.ico          <-- assets/favicon.ico landed at the ROOT
target/site/pkg
target/site/pkg/egonik-site.wasm
target/site/pkg/egonik-site.js
target/site/pkg/egonik-site.css

$ curl -s -o /dev/null -w '%{http_code}' localhost:3000/favicon.ico
200
\end{lstlisting}
\end{verified}

So a file you place at \file{assets/media/me.jpg} is served at \textbf{\code{/media/me.jpg}} --- not
at \code{/assets/media/me.jpg}. This matters concretely: \code{personal\_informations.image\_url}
is a column you are about to fill in, and the intuitive value is the wrong one.

\begin{pitfall}[And \code{/assets/} currently serves the entire site root]
\file{main.rs} contains \code{Files::new("/assets", \&site\_root)} --- note that the second argument
is the site root, not the assets directory. So the URL prefix \code{/assets} is mapped onto
everything.
\begin{lstlisting}[style=out]
$ curl -s -o /dev/null -w '%{http_code} %{size_download}' \
       localhost:3000/assets/pkg/egonik-site.wasm
200 2874275
\end{lstlisting}
Every file is reachable by two different URLs. Nothing under \file{target/site} is secret, so this
is not an information leak --- but two URLs for one resource is bad for caching and confusing to
debug. It is inherited from the upstream template. Defect \dnum{18}.
\end{pitfall}

\section{The check that proves nothing}

A closing warning for this part, because it undermines the habit recommended in
Chapter~\ref{ch:twobuilds} if you get it wrong.

\begin{verified}[Verified: a bare \code{cargo check} compiles neither program]
\begin{lstlisting}[style=out]
$ cargo check --no-default-features
    Finished `dev` profile [unoptimized + debuginfo] target(s)
\end{lstlisting}
It succeeds --- and it has checked \emph{neither} the server nor the browser build. With no features
enabled, \code{leptos} has no renderer, every server module is \code{cfg}'d out, and
\file{main.rs} falls through to its \code{\#[cfg(not(any(feature = "ssr", feature = "csr")))]} stub,
which is an empty \code{fn main()}.
\end{verified}

A green \code{cargo check} in this repository is therefore meaningless unless it names a feature.
Always both, always explicitly:

\begin{lstlisting}[style=sh]
cargo check --no-default-features --features ssr
cargo check --lib --no-default-features --features hydrate --target wasm32-unknown-unknown
\end{lstlisting}
